\documentclass[10pt,a4paper]{article}
\usepackage[margin=1in]{geometry}
\usepackage{hyperref}
\hypersetup{
    colorlinks=true,
    linkcolor=blue,
    urlcolor=blue,
    citecolor=blue
}
\usepackage{enumitem}
\usepackage{titlesec}
\usepackage{parskip}
\usepackage{fontawesome5}  

\titleformat{\section}{\large\bfseries}{}{0em}{}[\titlerule]

\begin{document}

\begin{center}
    {\Huge \textbf{Fernando Gutiérrez Canales}}\\
    {\Large Analista de Datos | Doctorado (PhD)}\\    
    León, México \quad | \quad \href{https://www.linkedin.com/in/fernando-guti%C3%A9rrez-canales-804876343/}{LinkedIn} \quad | \quad +52 1 477 55 18 832 \quad | \quad \href{mailto:carl.cfgc@gmail.com}{carl.cfgc@gmail.com} \\
     \href{https://github.com/Fernando-Canales}{GitHub: github.com/Fernando-Canales}
\end{center}

\vspace{1em}

\section*{Resumen Profesional}
Analista de Datos con Doctorado en Ciencias Exactas, reciente certificación en un Bootcamp de Análisis de Datos y más de 5 años de experiencia analizando conjuntos de datos complejos. Historial comprobado en la extracción, limpieza y transformación de datos sin procesar en inteligencia de negocios accionable utilizando SQL, Python (Pandas) y Tableau. Altamente motivado por la resolución de problemas, con gran habilidad para comunicar hallazgos técnicos a las partes interesadas (\textit{stakeholders}) y optimizar flujos de trabajo operativos para impulsar la toma de decisiones basada en datos.

\section*{Experiencia Profesional}
\textbf{Mercor} \hfill \textbf{Göttingen (Remoto)} \\
\textit{Ingeniero de Prompts de IA para Laboratorio de IA de Primer Nivel} \hfill \textit{Oct 2025 -- Feb 2026}
\begin{itemize}[leftmargin=1.5em]
	\item Desarrollé conjuntos de datos lógicos y analíticos complejos (prompts adversarios) para entrenar Modelos de Lenguaje Grande (LLMs), probando y mejorando efectivamente su capacidad de razonamiento en un 30\%.
\end{itemize}

\textbf{Observatorio de París \& Instituto Max Planck de Investigaciones del Sistema Solar} \hfill \textbf{París y Göttingen} \\
\textit{Analista de Datos e Investigador Doctoral} \hfill \textit{Mar 2022 -- Jul 2025}
\begin{itemize}[leftmargin=1.5em]
    \item Desarrollé flujos de trabajo automatizados (\textit{pipelines}) en Python para procesar y analizar volúmenes masivos de datos, extrayendo métricas clave para estimar eficiencias de detección de exoplanetas para la misión espacial PLATO.
    \item Colaboré con \textit{stakeholders} de más de 15 instituciones internacionales, estandarizando flujos de datos y asegurando un control riguroso de calidad e integridad de los datos para información crítica de la misión.
    \item Integré módulos de C y Bash para optimizar los procesos operativos, reduciendo significativamente el tiempo de manejo manual de datos.
\end{itemize}

\textbf{ESTEC, Agencia Espacial Europea (ESA)} \hfill \textbf{Noordwijk, Países Bajos}\\
\textit{Analista de Datos (Prácticas)}  \hfill \textit{Jun 2023 -- Ago 2023}
\begin{itemize}[leftmargin=1.5em]
    \item Automaticé la extracción y validación de parámetros complejos de detectores fotométricos utilizando Python, mejorando la eficiencia operativa y la precisión de los procesos de reporte de datos.
\end{itemize}

\section*{Proyectos Destacados de Datos}
\textbf{Proyecto de Analítica de Negocios / E-commerce (TripleTen)} \hfill \textbf{Python, SQL, Tableau} \\
\textit{Analista de Datos} \hfill \textit{Marzo 2026 -- Abril 2026}
\begin{itemize}[leftmargin=1.5em]
    \item Analicé datos de comportamiento de usuarios para identificar puntos de abandono (\textit{drop-offs}) en el embudo de ventas.
    \item Creé un \textit{dashboard} en Tableau que destacó una oportunidad del 15\% para aumentar la retención de clientes.
\end{itemize}

\textbf{Digitalización de Datos y BI para Clínica Oftalmológica} \hfill \textbf{León, México} \\
\textit{Consultor de Datos Principal (Proyecto Independiente)} \hfill \textit{2021 -- 2022}
\begin{itemize}[leftmargin=1.5em]
    \item Extraje, limpié y estructuré años de registros físicos y digitales fragmentados, centralizándolos en una base de datos robusta y de fácil consulta.
    \item Diseñé y generé reportes operativos recurrentes que mejoraron la accesibilidad a la información, optimizando directamente la gestión del personal y los flujos de trabajo diarios de la clínica.
\end{itemize}

\section*{Habilidades Técnicas}
\textbf{Análisis de Datos y BI:} SQL, Tableau, Excel, Visualización de Datos, Inteligencia de Negocios (BI), Limpieza de Datos (\textit{Data Wrangling}), Análisis de Series de Tiempo. \\
\textbf{Programación y Herramientas:} Python (Pandas, NumPy, Scikit-learn), Bash, C, Git, Docker, Linux. \\
\textbf{Habilidades Blandas:} Resolución de Problemas, Comunicación con Stakeholders, Pensamiento Analítico, Adaptabilidad.

\section*{Educación}
\textbf{Certificación de Bootcamp en Análisis de Datos} \hfill \textit{TripleTen} \hfill \textit{2026}\\
\textbf{Doctorado en Astrofísica} \hfill \textit{Observatorio de París \& MPS Göttingen} \hfill \textit{2025}\\
\textbf{Maestría en Astrofísica} \hfill \textit{Universidad de Guanajuato, México} \hfill \textit{2021}\\
\textbf{Licenciatura en Física} \hfill \textit{Universidad de Guanajuato, México} \hfill \textit{2019}

\section*{Logros y Publicaciones}
\begin{itemize}[leftmargin=1.5em]
    \item Coautor de \href{https://ui.adsabs.harvard.edu/search/fq...}{publicaciones revisadas por pares} sobre análisis de datos fotométricos y sistemas planetarios.
    \item Seleccionado para presentar hallazgos de investigación analítica en conferencias internacionales de primer nivel (EAS Padua 2024, PLATO Week Praga 2023).
    \item Beneficiario de becas académicas completas y graduado con Mención Honorífica en los tres grados universitarios (2019, 2021, 2025).
\end{itemize}

\end{document}
